\documentclass{article}
\usepackage{graphicx}
\usepackage{float}
\usepackage{caption}
\usepackage[en-GB]{datetime2}
\DTMlangsetup[en-GB]{ord=omit}
\graphicspath{{./images}}
\captionsetup{textfont=small, labelfont=bf}


\title{Applied Statistics - Assessment 1}
\author{Fernaldy Aristo Wirjowerdojo}
\date{\today}

\begin{document}

\maketitle

\section*{Task 1}


\pagebreak
\section*{Task 2}
\subsection*{a. Does Thyroid-stimulating hormone differ between male and female participants?}
Thyroid-stimulating hormone \textit{TSH} is a continuous variable while \textit{gender} is binary. So, the
independent t-test will be used and the assumptions for that test will be checked first. The following
hypothesis is considered:
\begin{itemize}
    \item $H_0$: There is no difference in the mean thyroid-stimulating hormone levels between male and female participants in the population
    \item $H_0$: There is a difference in the mean thyroid-stimulating hormone levels between male and female participants in the population
\end{itemize}

The assumptions to check are:
\begin{itemize}
    \item Independent groups
    \item TSH levels are normally distributed in both groups
    \item Variance of the TSH levels are equal between both groups
\end{itemize}

\begin{figure}[htbp]
    \centering
    \includegraphics[width=\textwidth]{2a-1.png}
    \caption{Distribution of TSH between male and female}
    \label{fig:2a-1}
\end{figure}

\begin{figure}[htbp]
    \centering
    \includegraphics[width=\textwidth]{2a-2.png}
    \caption{Levene's test for equality of variance of TSH}
    \label{fig:2a-2}
\end{figure}

All the assumptions are met, both groups are independent of one another, the TSH variable is normally distributed
in both groups, and the Levene's test ($p = 0.579 > 0.05$) indicates that the variance is homogeneous.

\begin{figure}[htbp]
    \centering
    \includegraphics[width=\textwidth]{2a-3.png}
    \caption{Independent t-test result}
    \label{fig:2a-3}
\end{figure}

Since the p-value is significant ($p < 0.001$), the $H_0$ is rejected. There is a difference in the mean
thyroid-stimulating hormone levels between male and female participants in the population.
\\
\textbf{In conclusion}, female participants had lower thyroid-stimulating hormone levels with an average difference
of $0.886 \ mg/dL$. Furthermore, the 95\% confidence interval ($-0.899, -0.873$) indicates that mean
thyroid-stimulating hormone levels in the female population is between 0.873 and 0.899 lower than that of the
male population.


% 2b. Kruskal-Wallis
\pagebreak
\subsection*{b. Does the level of Triglycerides differ among different levels of physical activity?}
Triglycerides \textit{triglycerides} is a continuous variable while physical activity
\textit{physical\_activity} is a categorical variable. The hypothesis for this question are the following:
\begin{itemize}
    \item $H_0$: There is no difference in the distribution of triglycerides across all physical activity groups in the population
    \item $H_1$: There is a difference in the distribution of triglycerides across all physical activity groups in the population
\end{itemize}
Initially, the first test to consider is ANOVA. So the following three assumptions are checked.
\begin{itemize}
    \item The three physical activity groups are independent
    \item The distribution of \textit{triglycerides} is normal in the three groups
    \item The variance of \textit{triglycerides} is equal between the three groups
\end{itemize}

\begin{figure}[htbp]
    \centering
    \includegraphics[width=\textwidth]{2b-1.png}
    \caption{Distribution of triglycerides between the three physical activity groups}
    \label{fig:2b-1}
\end{figure}

The first assumption is met but the figure above shows that the second assumption fails because the
distribution of triglycerides for all the three groups are not normal and instead are right-skewed.
Instead of the regular ANOVA, the Kruskal-Wallis test is used and the hypothesis is adjusted for non-parametric.
Before the test results, the median and the IQR can be considered.

\begin{table}[H]
\centering
\begin{tabular}{|l|l|l|l|}
\hline
\textbf{Activity Levels} & \textbf{Median} & \textbf{Lower Quartile} & \textbf{Upper Quartile} \\ \hline
\textbf{No Activity}        & 153.50 & 104.00 & 228.00 \\ \hline
\textbf{Moderate Activity}  & 137.50 & 97.00  & 192.00 \\ \hline
\textbf{Strenuous Activity} & 130.00 & 97.00  & 178.00 \\ \hline
\end{tabular}
\caption{Median and IQR of triglycerides levels across physical activity}
\label{tab:2b-descriptive-stats}
\end{table}

The median from the three groups show that the triglycerides levels are generally lower as the level of
physical activity gets more strenuous. But the Kruskal-Wallis test will be used to determine its significance.

\begin{figure}[htbp]
    \centering
    \includegraphics[width=\textwidth]{2b-2.png}
    \caption{Kruskal-Wallis test result}
    \label{fig:2b-2}
\end{figure}
The test result shows that the difference is significant ($p < 0.001$). The null hypothesis is rejected.
\\~\\
\textbf{In conclusion}, the distribution of triglycerides is different across the physical activity groups of no activity,
moderate activity, and strenuous activity in the population. The triglycerides levels are lower as more intense
the physical activity is.





\pagebreak
\subsection*{c. Is there an association between physical activity and gender?}


% 2d. ANOVA
\pagebreak
\subsection*{d. Does low-density lipoprotein differ among different BMI categories?}
Low-density lipoprotein \textit{LDL} is a continuous variable while BMI categories \textit{BMI\_cat} is a categorical
variable. The hypothesis for this question are the following:
\begin{itemize}
    \item $H_0$: There is no difference in the mean of low-density lipoprotein levels between the BMI categories: under \& normal weight, overweight, and obese in the population
    \item $H_1$: There is a difference in the mean of low-density lipoprotein levels between the BMI categories: under \& normal weight, overweight, and obese in the population
\end{itemize}

Since this is a continuous against a categorical variable, ANOVA will be used to determine the significance. But,
the assumptions will be checked first for ANOVA:
\begin{itemize}
    \item The BMI categories are independent
    \item The distribution of ldl is normal for the three BMI categories
    \item The variance of ldl is equal for the three BMI categories
\end{itemize}

\begin{figure}[htbp]
    \centering
    \includegraphics[width=\textwidth]{2d-1.png}
    \caption{Distribution of ldl between the three BMI categories}
    \label{fig:2d-1}
\end{figure}

\begin{figure}[htbp]
    \centering
    \includegraphics[width=\textwidth]{2d-2.png}
    \caption{Levene's test for equality of variance of ldl}
    \label{fig:2d-2}
\end{figure}

The first assumption is immediately met. Figure 4 shows that the distribution of LDL across the three BMI categories
are normal and Figure 5 for the result of Levene's test of homogeneity of variance $(p = 0.228 > 0.05)$ indicates
that the variance is equal across the three BMI categories. All three assumptions for ANOVA is met.


\begin{figure}[H]
    \centering
    \includegraphics[width=\textwidth]{2d-3.png}
    \caption{ANOVA results}
    \label{fig:2d-3}
\end{figure}

\begin{figure}[H]
    \centering
    \includegraphics[width=\textwidth]{2d-4.png}
    \caption{Post-hoc Scheffe test}
    \label{fig:2d-4}
\end{figure}


\textbf{In conclusion}, The ANOVA test $(p < 0.001)$ indicates that the $H_0$ is rejected.
Thus, there is a difference in the mean of low-density lipoprotein across
the three BMI categories under \& normal weight, overweight, and obese in the population. Furthermore, the
post-hoc test shows the significant difference in the mean between the groups:
\begin{itemize}
    \item Under \& normal weight and overweight is around $19.3$
    \item Under \& normal weight and obese is around $19.8$
\end{itemize}
While the difference between the overweight and obese groups is not statistically significant.


\pagebreak
\subsection*{e. Is there an association between smoking and CAD?}


\pagebreak
\subsection*{f. Is there an association between congenital anomalies and CAD?}


% 2g. Paired t-test
\pagebreak
\subsection*{g. Is there a change in Total cholesterol between baseline and after intervention?}
Total cholesterol base and after (\textit{Tchol\_base} and \textit{Tchol\_after}) are continuous variables
and since it wants to know if the change between baseline and after, the paired t-test will be considered and
the assumption \textit{The difference between total cholesterol at baseline and after intervention
is normally distributed} will be checked first and the following hypothesis will be used:
\begin{itemize}
    \item $H_0$: The mean total cholesterol between baseline and after intervention is the same in the population
    \item $H_1$: The mean total cholesterol between baseline and after intervention is not the same in the population
\end{itemize}

The new variable \textit{diff\_tchol} is introduced by subtracting \textit{Tchol\_base} from \textit{Tchol\_after}
\begin{figure}[H]
    \centering
    \includegraphics[width=\textwidth]{2g-1.png}
    \caption{Difference in total cholesterol at baseline and after intervention}
    \label{fig:2g-1}
\end{figure}
The figure above shows the new variable is normally distributed. So, the assumption is met.

\begin{figure}[H]
    \centering
    \includegraphics[width=\textwidth]{2g-2.png}
    \caption{Paired t-test}
    \label{fig:2g-2}
\end{figure}

The t-test result ($p < 0.001$) shows that it is significant and that the $H_0$ should be rejected.

\textbf{In conclusion}, the mean total cholesterol between baseline and after intervention is
different in the population. The average total cholesterol after intervention is around $40 \ mg/dL$
lower. Whereas the 95\% confidence interval for the mean difference is ($39.97, 40.85$) indicating
that the mean total cholesterol after intervention in the population is between $39.97 \ mg/dL$ lower
or $40.85 \ mg/dL$ higher.


\end{document}
